\subsection{La comparación entrópica en el tramo macroscópico}

\begin{lema}[Factorización estacionaria]\label{coste-vacancias}
Sea $p$ un punto estacionario interior de la entropía de Bethe $H_B$, con masas de fila y de columna
iguales a uno y $h_f>0$. Existen números positivos $a_i,b_j$ tales que, en toda casilla disponible,
\[
 v_e:=p_e(1-p_e)^3=a_i b_j h_{D(e)}h_{A(e)} .
\]
Pongamos $\mathcal K=\prod_i a_i\prod_j b_j\prod_f h_f$ y $W=\sum_e v_e$. Entonces
$H_B(p)=-\log\mathcal K+3\sum_e\log(1-p_e)$. Si además $p_{\max}\le C/k$, para $k$
suficientemente grande se tiene $W=k+O_C(1)$ y $H_B(p)=-\log\mathcal K-3k+O_C(1)$.
\end{lema}
\begin{proof}
La derivada de $H_B$ con respecto a $p_e$ es
$-\log p_e-3\log(1-p_e)+\log h_{D(e)}+\log h_{A(e)}-2$.
En el punto interior se anula su producto escalar con toda variación que preserve las sumas de filas y
columnas. Por ortogonalidad al núcleo de la aplicación que da esas sumas, esta derivada es una suma de un
multiplicador de fila y otro de columna. Al exponenciar y absorber la constante se obtiene la
factorización. Multiplicando su logaritmo por $p_e$ y sumando en las casillas, los dos multiplicadores
contribuyen $\sum_i\log a_i+\sum_j\log b_j$ y las diagonales contribuyen $\sum_f\ell_f\log h_f$. La
identidad para $H_B$ se sigue de $\ell_f+h_f=1$. Además $0\le p_e-v_e\le3p_e^2$ y
$|\log(1-p_e)+p_e|\le p_e^2/(1-p_{\max})$. Ambas cotas se pueden sumar, pues $\sum_e p_e^2\le C$.
\end{proof}

\begin{lema}[Comparación con colocaciones aumentadas]\label{coste-comparacion}
Sea $B$ un tablero con $Q(B)>0$ y con un punto estacionario interior de Bethe $p$ de masas de fila y
de columna iguales a uno y vacancias diagonales positivas. Usamos los pesos $v_e=p_e(1-p_e)^3$ de las
piezas reales y el producto $\mathcal K$ de la factorización del Lema~\ref{coste-vacancias}. Si $n$
es el orden del tablero inicial, pongamos $N=2n-1$ y $H=N-k$.
Existe una ley de comparación, uniforme sobre colocaciones aumentadas ordenadas, con $k$ piezas
reales y $H$ piezas
virtuales, para la cual
\[
 \log Q(B)\le
 \E\sum_{s=1}^N\log W_{T_s}(s)
       -\log\mathcal K-\log k!-4\log H! .
\]
Aquí $T_s$ es el tipo de la pieza en el lugar $s$, y $W_{T_s}(s)$ es la masa disponible de ese tipo
inmediatamente antes de revelar la pieza.
\end{lema}
\begin{proof}
Desarrollamos la cuenta de colocaciones aumentadas de \cite[Teorema~3.6]{ultimas} para el soporte de
$B$. En cada una de las cuatro familias tomamos $N$ vértices. En las dos primeras conservamos las $k$
filas y las $k$ columnas físicas de $B$ y añadimos $H$ vértices ficticios en cada familia. En cada
familia diagonal usamos las $N$ etiquetas del tablero inicial. Una diagonal residual conserva su línea
física. En el lugar de una diagonal ya ocupada ponemos un vértice ficticio de la misma etiqueta, al
que ninguna casilla residual es incidente. Una línea sin casillas disponibles tiene también grado real
cero. En ambos casos la masa residual es cero y extendemos $h_f=1$ a esa etiqueta. Esta extensión sólo
interviene en la cuenta auxiliar. Deja inalteradas la identidad de estacionariedad en todas las
casillas disponibles y la constante $\mathcal K$.

Las piezas reales son las casillas disponibles con peso $v_e$. Para las piezas virtuales asignamos
peso de vértice cero a filas y columnas físicas, uno a filas y columnas ficticias y $h_f$ a las
etiquetas diagonales. Una pieza virtual consta de un vértice en cada familia y pesa el producto de sus
cuatro pesos. Sólo usamos las de peso positivo. En una familia diagonal cada casilla contribuye su
masa exactamente una vez. Por ello
\[
 \sum_{f\text{ de una familia diagonal}}h_f=N-\sum_e p_e=H.
\]
Los pesos de vértice suman igualmente $H$ en las otras dos familias. De aquí que la masa virtual
inicial sea $H^4$, aunque hubiese diagonales ocupadas antes de llegar a $B$.

Una colocación real completa deja descubiertos $H$ vértices en cada familia. Para completarla fijamos
las $H$ filas ficticias como índices y elegimos independientemente tres biyecciones con los vértices
descubiertos de las otras tres familias. Esto da exactamente $(H!)^3$ completaciones virtuales, todas
de peso positivo porque cada $h_f>0$. Recíprocamente, en una colocación aumentada que cubre todos los
vértices, las filas físicas deben estar cubiertas por piezas reales y las ficticias por virtuales. Las
piezas reales forman así una de las $Q(B)$ completaciones de $B$. Cada colocación aumentada admite $N!$
órdenes. En cualquiera de ellos el producto de los pesos de todas las piezas es $\mathcal K$. En
efecto, las reales aportan una vez cada $a_i$ y cada $b_j$, y entre reales y virtuales se usa una vez
el factor $h_f$ de cada etiqueta diagonal.

Como en~\cite[Sección~3]{ultimas}, sea $\mathbb P$ la ley de comparación, uniforme sobre estas
historias exitosas, y sea $\mathbb Q$ el siguiente proceso de referencia. Primero se toma
uniformemente una palabra de colores con $k$ colores reales y $H$ virtuales. En cada
lugar se elige entre las piezas disponibles del tipo prescrito con probabilidad proporcional a su
peso. Una pieza está disponible si ninguno de sus vértices ha sido cubierto. Si la masa de ese tipo es
cero, el proceso termina en fallo y conserva allí la probabilidad correspondiente. Para una historia
exitosa $\omega$ se tiene
\[
 \mathbb P(\omega)=\frac{1}{Q(B)(H!)^3N!},\qquad
 \mathbb Q(\omega)=\frac{k!H!}{N!}\,
             \frac{\mathcal K}{\prod_{s=1}^N W_{T_s}(s)} .
\]
En particular $\mathbb Q(\omega)>0$. Jensen aplicado al logaritmo da
$\E_{\mathbb P}\log(\mathbb Q/\mathbb P)\le\log\sum_{\omega:\mathbb P(\omega)>0}\mathbb Q(\omega)\le0$.
Sustituyendo ambas masas se obtiene la desigualdad del lema. La esperanza se toma bajo $\mathbb P$,
en el
único tablero $B$ que se está contando.
\end{proof}
